In this section, we give the proof of our main theorem, \Cref{thm: d4 hj3}, modulo three inputs from the later sections of the paper. Our core strategy is to use the motivic zig-zag to gain information about the classical Adams $d_4$ differential on $h_j^3$ from the associated algebraic Novikov $d_4$ differential.

%For indexes, recall from Section~\ref{sec:review} that, for elements in the classical Adams $E_2$-page $\Ext_{A}^{a,t}(\mathbb{F}_2, \mathbb{F}_2)$, we have $a$ is the Adams filtration, $t$ is the internal degree; for elements in the motivic Adams $E_2$-page $\Ext^{a,t,w}_{A^\textup{mot}}(\mathbb{F}_2[\tau], \mathbb{F}_2[\tau])$, we have $w$ is the motivic filtration. Moreover, from the motivic Cartan-Eilenberg spectral sequence, every element in $\Ext^{a,t,w}_{A^\textup{mot}}(\mathbb{F}_2[\tau], \mathbb{F}_2[\tau])$ has a Cartan-Eilenberg filtration, denoted by $k$, and an Adams-Novikov filtration $s$, satisfying that $a = s+k$.

\begin{rec}
  As in the classical Steenrod algebra, $\xi_1^{2^j}$ is a primitive element in the motivic dual Steenrod algebra. We let $h_j$ denote the associated class on the $E_2$-page of the motivic Adams spectral sequence. $h_j$ lives in $(a,s,t,w)$-degree $(1, \ 1, \ 2^j, \ 2^{j-1})$. Similarly, we also denote the induced classes on the $E_2$-pages of the classical Adams sseq and motivic Adams sseq for $S^{0,0}/\tau$ by $h_j$.
\end{rec}

In fact, the main result of this section is a computation of $d_4(h_j^3)$ in the motivic Adams spectral sequence (with \Cref{thm: d4 hj3} being obtained from this by Betti realization). Before proceeding let us summarize what we need from later sections.

\begin{itemize}
\item In \Cref{E2 descriptions} we describe the $E_2$-pages of all three spectral sequences in the motivic zig-zag (classical Adams, motivic Adams, motivic Adams for $S^{0,0}/\tau$) in a neighborhood of the classes $h_j^3$. In particular, the Adams filtration 6 and 7 lines in \Cref{E2 descriptions}~(1) are the crucial new information. This theorem is proved in Sections \ref{sec:algAH} and \ref{sec:cess} and summarized in \Cref{fig:3charts}.
\item In \Cref{diff motivic ctau} we compute the differentials on $h_j^3$ in the motivic Adams sseq for $S^{0,0}/\tau$. This theorem is proved in \Cref{sec:alg-nov-diff}.
\item In \Cref{prop:diff classical} we give partial information on the classical Adams differentials on $h_j^3$.
  This theorem is proved in \Cref{sec:kervaire} using $\F_2$-synthetic homotopy theory.
\end{itemize}

% \begin{thm} \label{E2 descriptions}
% %Near the elements $h_j^3$ for $j \geq 6$, we have the following description of the $E_2$-pages of the classical Adams spectral sequence for $S^{0}$ and the motivic Adams spectral sequences for $S^{0,0}$ and $S^{0,0}/\tau$.  (See Figure.)
% Suppose that $j \geq 6$.
% \begin{enumerate}
% \item In the $E_2$-page of the classical Adams spectral sequence for $S^{0}$, 
% $$E_2 = \Ext_{A}^{a,t}(\mathbb{F}_2, \mathbb{F}_2) \Rightarrow \pi_{t-a}S^{0},$$
% \begin{enumerate}
% \item in the range
% $$a \leq 5, \ t-a = 3\cdot 2^j - 3, \ 3\cdot 2^j - 4$$
% the only nonzero elements are 
% $$h_j^3, \ h_0 h_j^3, \ h_0^2 h_j^3, \ g_{j-2}, \ h_0g_{j-2}, \ h_1g_{j-2};$$
% \item in the range
% $$a = 6, \ 7, \ t-a = 3\cdot 2^j - 3,$$ 
% the following elements are nonzero:
% $h_0^2g_{j-2}, \ h_0^3g_{j-2}$.
% \end{enumerate}

% \item In the $E_2$-page of the motivic Adams spectral sequence for $S^{0,0}$, 
% $$E_2 = \Ext^{a,t,w}_{A^\textup{mot}}(\mathbb{F}_2[\tau], \mathbb{F}_2[\tau]) \Rightarrow \pi_{t-a, w}S^{0,0},$$
% $$\Ext^{s,t,w}_{P^\textup{mot}}(\F_2[\tau], \Ext^k_{Q^\textup{mot}}(\F_2[\tau], \F_2[\tau])) \Rightarrow \Ext^{s+k, t, w}_{A^\textup{mot}}(\F_2[\tau], \F_2[\tau]),$$
% \begin{enumerate}
% \item in the range
% $$a \leq 4, \ t-a = 3\cdot 2^j - 3, \ 3\cdot 2^j - 4,$$
% the only nonzero groups are copies of $\F_2[\tau]$, generated by 
% $$h_j^3, \ h_0 h_j^3, \ g_{j-2},$$
% all with the same weight $w = 3 \cdot 2^{j-1}$,\\
% and possible copies of $\F_2$, generated by elements with degree
% $$(a, \ t-a, \ s, \ w) = (4, \ 3\cdot 2^j - 3, \ 3, \ 3\cdot 2^{j-1}).$$

% \item In the range
% $$a = 5, \ t-a = 3\cdot 2^j - 3, \ 3\cdot 2^j - 4,$$
% the only nonzero groups are copies of $\F_2[\tau]$, generated by 
% $$h_0^2 h_j^3, \ h_0g_{j-2}, \ h_1g_{j-2},$$
% possible copies of $\F_2[\tau]/\tau^2$, generated by elements with degree
% $$(t-a, \ s,\ w)=(3\cdot 2^j - 3, \ 5,\ 3\cdot 2^{j-1}+1),$$
% and possible copies of $\F_2$, generated by elements with degrees 
% $$(t-a, \ s, \ w) = (3\cdot 2^j - 3, \ 5, \ 3\cdot 2^{j-1} +1), \ (3\cdot 2^j - 3, \ 3, \ 3\cdot 2^{j-1}), \ (3\cdot 2^j - 4, \ 4, \ 3\cdot 2^{j-1}).$$

% \item In the range
% $$a = 6, \ t-a = 3\cdot 2^j - 4,$$
% the only nonzero groups are copies of $\F_2[\tau]$, generated by $h_0^2g_{j-2}$, and\\ 
% possibly other elements with degrees
% $$ (s,\ w)=(0, \ 3\cdot 2^{j-1}-2), \ (2, \ 3\cdot 2^{j-1}-1), \ (4,\ 3\cdot 2^{j-1}),\ (6,\ 3\cdot 2^{j-1}+1),$$
% %possibly copies of $\F_2[\tau]/\tau^2$, generated by elements with degree
% %$$(s,\ w)=(6,\ 3\cdot 2^{j-1}+1),$$
% and possible copies of $\F_2$, generated by elements with degrees
% $$(s,\ w)=(4,\ 3\cdot 2^{j-1}),\ (6,\ 3\cdot 2^{j-1}+1).$$
% In particular, there are $\textbf{no copies}$ of $\F_2[\tau]/\tau^2$, generated by elements with degree
% $$(s,\ w)=(6,\ 3\cdot 2^{j-1}+1).$$


% \item In the range
% $$a = 7, \ t-a = 3\cdot 2^j - 4,$$
% the only nonzero groups are copies of $\F_2[\tau]$, generated by $h_0^3g_{j-2}$, and\\ 
% possibly other elements with degrees
% $$ (s,\ w)=(0, \ 3\cdot 2^{j-1}-2), \ (2, \ 3\cdot 2^{j-1}-1), \ (4,\ 3\cdot 2^{j-1}),\ (6,\ 3\cdot 2^{j-1}+1),$$
% possible copies of $\F_2[\tau]/\tau^2$, generated by elements with degree
% $$(s,\ w)=(6,\ 3\cdot 2^{j-1}+1),$$
% and possible copies of $\F_2$, generated by elements with degrees
% $$(s,\ w)=(4,\ 3\cdot 2^{j-1}),\ (6,\ 3\cdot 2^{j-1}+1).$$
% \end{enumerate}

% \item In the the $E_2$-page of the motivic Adams spectral sequence for $S^{0,0}/\tau$, 
% $$E_2 = \Ext^{a,t,w}_{A^\textup{mot}}(\mathbb{F}_2[\tau], \mathbb{F}_2) \Rightarrow \pi_{t-a, w}S^{0,0}/\tau,$$
% in the range
% $$a \leq 7, \ t-a = 3\cdot 2^j - 3, \ 3\cdot 2^j - 4$$
% the following elements are nonzero:
% $$h_j^3, \ h_0 h_j^3, \ h_0^2 h_j^3, \ g_{j-2}, \ h_0g_{j-2}, \ h_0^2g_{j-2}, \ h_0^3g_{j-2}, \ h_1g_{j-2}.$$
% \end{enumerate}
% \end{thm}

\begin{thm} \label{E2 descriptions}
  Let $j \geq 6$.
  \begin{enumerate}
  \item The $E_2$-page of the classical Adams spectral sequence for $S^{0}$, 
    $$ \Ext_{\A}^{a,t}(\mathbb{F}_2, \mathbb{F}_2) \cong E_2^{a,t} \Rightarrow \pi_{t-a}S^{0},$$
    takes the following form near $h_j^3$:
    \begin{center}{\renewcommand{\arraystretch}{1.2}
        \begin{tabular}{|ll|c|}\hline
          $(a,$ & $t-a)$ & $\Ext^{a,\,t}_{\A}(\mathbb{F}_2,\, \mathbb{F}_2)$ \\\hline\hline
          $(4,$ & $ 3 \cdot 2^j-4)$ & $\F_2\{g_{j-2}\}$ \\\hline
          $(5,$ & $ 3 \cdot 2^j-4)$ & $\F_2\{h_0g_{j-2}\}$ \\\hline
          $(6,$ & $ 3 \cdot 2^j-4)$ & contains $\F_2\{h_0^2g_{j-2}\}$ \\\hline
          $(7,$ & $ 3 \cdot 2^j-4)$ & contains $\F_2\{h_0^3g_{j-2}\}$ \\\hline
          $(3,$ & $ 3 \cdot 2^j-3)$ & $\F_2\{h_j^3\}$ \\\hline
          $(4,$ & $ 3 \cdot 2^j-3)$ & $\F_2\{h_0h_j^3\}$ \\\hline
          $(5,$ & $ 3 \cdot 2^j-3)$ & $\F_2\{h_0^2h_j^3, h_1g_{j-2}\}$ \\\hline    
        \end{tabular}}
    \end{center}
    with the exception that in the case $j=6$ degree $(a,t-s)=(5,3 \cdot 2^{6} - 4)$ contains the additional class $h_7D_3(0)$, and that in the case $j=7$ degree $(a,t-s)=(5,3 \cdot 2^{7} - 3)$ contains the additional class $h_8D_3(1)$.
    
  \item The $E_2$-page of the motivic Adams spectral sequence for $S^{0,0}$, 
    $$\Ext^{a,t,w}_{\A^\textup{mot}}(\mathbb{F}_2[\tau], \mathbb{F}_2[\tau]) \cong E_2^{a,t,w} \Rightarrow \pi_{t-a, w}S^{0,0},$$
    takes the following form near $h_j^3$:
    \begin{center}{\renewcommand{\arraystretch}{1.2}
        \begin{tabular}{|ll|c|}\hline
          $(a,$ & $t-a)$ & $\Ext^{a,\,t,\,w}_{\A^\textup{mot}}(\mathbb{F}_2[\tau],\, \mathbb{F}_2[\tau])$ \\\hline\hline
          $(4,$ & $ 3 \cdot 2^j-4)$ & $\F_2[\tau]\{g_{j-2}\}$ \\\hline
          $(5,$ & $ 3 \cdot 2^j-4)$ & $\F_2[\tau]\{h_0g_{j-2}\}$ \\\hline
          $(6,$ & $ 3 \cdot 2^j-4)$ & contains $\F_2[\tau]\{h_0^2g_{j-2}\}$, $\tau$-torsion free \\\hline
          $(7,$ & $ 3 \cdot 2^j-4)$ & contains $\F_2[\tau]\{h_0^3g_{j-2}\}$, $\tau$-torsion free \\\hline
          $(3,$ & $ 3 \cdot 2^j-3)$ & $\F_2[\tau]\{h_j^3\}$ \\\hline
          $(4,$ & $ 3 \cdot 2^j-3)$ & $\F_2[\tau]\{h_0h_j^3\}$ \\\hline
          $(5,$ & $ 3 \cdot 2^j-3)$ & $\F_2[\tau]\{h_0^2h_j^3, h_1g_{j-2}\}$ \\\hline    
        \end{tabular}}
    \end{center}
    where each of the generators
    $g_{j-2}, h_0g_{j-2}, h_0^2g_{j-2}, h_0^3g_{j-2}, h_j^3, h_0h_j^3, h_0^2h_j^3$
    has weight $w = 3 \cdot 2^{j-1}$ and $h_1g_{j-2}$ has weight $w = 3 \cdot 2^{j-1}+1$, with the exception that
    in the case $j=6$ degree $(a,t-a)=(5,3 \cdot 2^{6} - 4)$ contains an additional $\F_2[\tau]$ summand with generator $h_7D_3(0)$, and that
    in the case $j=7$ degree $(a,t-a)=(5,3 \cdot 2^{7} - 3)$ contains an additional $\F_2[\tau]$ summand with generator $h_8D_3(1)$ of weight $w = 3 \cdot 2^{6}+1$.
    
  \item The $E_2$-page of the motivic Adams spectral sequence for $S^{0,0}/\tau$, 
    $$\Ext^{a,t,w}_{\A^\textup{mot}}(\mathbb{F}_2[\tau], \mathbb{F}_2) \cong E_2^{a,t,w} \Rightarrow \pi_{t-a, w}S^{0,0}/\tau,$$
    takes the following form near $h_j^3$:
    \begin{center}{\renewcommand{\arraystretch}{1.2}
        \begin{tabular}{|lll|c|}\hline
          $(a,$ & $2w-t+a,$ & $t-a)$ & $\Ext^{a,\,t,\,w}_{\A^\textup{mot}}(\mathbb{F}_2[\tau],\, \mathbb{F}_2)$ \\\hline\hline
          $(4,$ & $4,$ & $ 3 \cdot 2^j-4)$ & $\F_2\{g_{j-2}\}$ \\\hline
          $(5,$ & $4,$ & $ 3 \cdot 2^j-4)$ & $\F_2\{h_0g_{j-2}\} $ \\\hline
          $(6,$ & $4,$ & $ 3 \cdot 2^j-4)$ & contains $\F_2\{h_0^2g_{j-2}\} $ \\\hline
          $(7,$ & $4,$ & $ 3 \cdot 2^j-4)$ & contains $\F_2\{h_0^3g_{j-2}\} $ \\\hline
          $(3,$ & $1,$ & $ 3 \cdot 2^j-3)$ & $0$ \\\hline
          $(4,$ & $1,$ & $ 3 \cdot 2^j-3)$ & $0$ \\\hline
          $(5,$ & $1,$ & $ 3 \cdot 2^j-3)$ & $0$ \\\hline
          $(6,$ & $1,$ & $ 3 \cdot 2^j-3)$ & $0$ \\\hline          
          $(3,$ & $3,$ & $ 3 \cdot 2^j-3)$ & $\F_2\{h_j^3\}$ \\\hline
          $(4,$ & $3,$ & $ 3 \cdot 2^j-3)$ & $\F_2\{h_0h_j^3\}$ \\\hline
          $(5,$ & $3,$ & $ 3 \cdot 2^j-3)$ & $\F_2\{h_0^2h_j^3\}$ \\\hline    
          $(6,$ & $3,$ & $ 3 \cdot 2^j-3)$ & $\F_2\{h_0^3h_j^3\}$ or $0$ \\\hline
          $(5,$ & $5,$ & $ 3 \cdot 2^j-3)$ & $\F_2\{h_1g_{j-2}\}$ \\\hline          
        \end{tabular}}
    \end{center}
    where the generators have the same weights as in (2), with the exception that
    in the case $j=6$ degree $(a,2w-t+a,t-a)=(5,4,3 \cdot 2^{6} - 4)$ possibly contains an additional class $h_7D_3(0)$, and that
    in the case $j=7$ degree $(a,2w-t+a,t-a)=(5,5,3 \cdot 2^{7} - 3)$ contains an additional class $h_8D_3(1)$.

  \item Under the Betti realization and reduction mod $\tau$ maps each generator maps to the generator of the same name.
  \end{enumerate}
\end{thm}


\begin{sseqdata}[ name = 3charts1hj3, xscale=1.5, yscale=1.5, x range = {-4}{-3}, y range = {2}{7}, x tick step = 1, y tick step = 1, grid = crossword, Adams grading, lax degree]
  
  \class[circle, inner sep=0.5pt, "?", {font = \tiny}](-4,6)
  \class[circle, inner sep=0.5pt, "?", {font = \tiny}](-4,7)
  
  \class[fill, "g_{j-2}" {below} ](-4,4)
  \class[fill](-4,5) \structline
  \class[fill](-4,6) \structline
  \class[fill](-4,7) \structline
  
  \class[fill](-3,5) \structline(-4,4)
  
  \class[fill, "h_j^3" {below} ](-3,3)
  \class[fill](-3,4) \structline
  \class[fill](-3,5) \structline
  
\end{sseqdata}


% \begin{sseqdata}[ name = 3charts2hj3, xscale=1.5, yscale=1.5, x range = {-4}{-3}, y range = {2}{7}, x tick step = 1, y tick step = 1, grid = crossword, Adams grading, lax degree]

% \class[rectangle, inner sep=3pt](-4,6)
%  \class[rectangle, inner sep=3pt](-4,7)
 
%  \class[white ,circlen=2](-4,6)
%   \class[circle, inner sep=0.5pt, "4", {font = \tiny}](-4,6)
%   \class[circle, inner sep=0.5pt, "6", {font = \tiny}](-4,6)
    
%   \class[inner sep=0.5pt, "6", {font = \tiny},circlen=2](-4,7)
%   \class[circle, inner sep=0.5pt, "4", {font = \tiny}](-4,7)
%   \class[circle, inner sep=0.5pt, "6", {font = \tiny}](-4,7)
  


%   \class[rectangle, fill, inner sep=3pt, "g_{j-2}" {below} ](-4,4)
%   \class[rectangle, fill, inner sep=3pt](-4,5) \structline
%   \class[rectangle, fill, inner sep=3pt](-4,6) \structline
%   \class[rectangle, fill, inner sep=3pt](-4,7) \structline
  
%   \class[rectangle, fill, inner sep=3pt](-3,5) \structline(-4,4)
  
%   \class[circle, inner sep=0.5pt, "3", {font = \tiny}](-3,4)
%   \class[circle, inner sep=0.5pt, "4", {font = \tiny}](-4,5)
  
%   \class[inner sep=0.5pt, "5", {font = \tiny}, circlen=2](-3,5)
%   \class[circle, inner sep=0.5pt, "3", {font = \tiny}](-3,5)
%   \class[circle, inner sep=0.5pt, "5", {font = \tiny}](-3,5)

  
%    \class[rectangle, fill, inner sep=3pt, "h_j^3" {below} ](-3,3)
%   \class[rectangle, fill, inner sep=3pt](-3,4) \structline
%   \class[rectangle, fill, inner sep=3pt](-3,5) \structline
      
% \end{sseqdata}

\begin{sseqdata}[ name = 3charts2hj3, xscale=1.5, yscale=1.5, x range = {-4}{-3}, y range = {2}{7}, x tick step = 1, y tick step = 1, grid = crossword, Adams grading, lax degree]

  \class[rectangle, inner sep=1.5pt, "?", {font = \tiny}](-4,6)
  \class[rectangle, inner sep=1.5pt, "?", {font = \tiny}](-4,7)
       
  \class[rectangle, fill, inner sep=3pt, "g_{j-2}" {below} ](-4,4)
  \class[rectangle, fill, inner sep=3pt](-4,5) \structline
  \class[rectangle, fill, inner sep=3pt](-4,6) \structline
  \class[rectangle, fill, inner sep=3pt](-4,7) \structline
  
  \class[rectangle, fill, inner sep=3pt](-3,5) \structline(-4,4)
    
  \class[rectangle, fill, inner sep=3pt, "h_j^3" {below} ](-3,3)
  \class[rectangle, fill, inner sep=3pt](-3,4) \structline
  \class[rectangle, fill, inner sep=3pt](-3,5) \structline
  
\end{sseqdata}


\begin{sseqdata}[ name = 3charts3hj3, xscale=1.5, yscale=1.5, x range = {-4}{-3}, y range = {2}{7}, x tick step = 1, y tick step = 1, grid = crossword, Adams grading, lax degree]

  \class[circle, inner sep=0.5pt, "?", {font = \tiny}](-4,6)
  \class[circle, inner sep=0.5pt, "?", {font = \tiny}](-4,7)
  
  \class[fill, "g_{j-2}" {below} ](-4,4)
  \class[fill](-4,5) \structline
  \class[fill](-4,6) \structline
  \class[fill](-4,7) \structline
  
  \class[fill](-3,5) \structline(-4,4)

  \class[circle, inner sep=0.5pt, "5", {font = \tiny}](-3,6)
  
   \class[fill, "h_j^3" {below} ](-3,3)
  \class[fill](-3,4) \structline
  \class[fill](-3,5) \structline
  \class[fill, red](-3,6) \structline
    
\end{sseqdata}




\begin{figure}[h]
  \centering
  The classical and motivic Adams spectral sequences \\
  for $S^0, \ S^{0,0}$ and $S^{0,0}/\tau$ near $h_j^3$
   %\\\vspace{6pt}
  \scalebox{1.00}{
    \printpage[ name = 3charts1hj3, page = 3, y axis gap = 0.8cm, right clip padding = 0cm, x axis extend end = 0.75cm ] \quad
    \printpage[ name = 3charts2hj3, page = 3, no y ticks, x axis tail = 0cm, y axis gap = 0.8cm, right clip padding = 0cm, x axis extend end = 0.75cm ] \quad
    \printpage[ name = 3charts3hj3, page = 3, no y ticks, x axis tail = 0cm, y axis gap = 0.8cm, right clip padding = 0cm, x axis extend end = 0.75cm ]
  }
  \caption{ 
    The horizontal degree is the topological stem $t-a$, shifted by $3\cdot 2^j$.     
    The vertical degree is the Adams filtration $a$. 
    Left and Right:
    Each dot $\bullet$ denotes a copy of $\F_2$. 
    Circles $\circ$ denote possible copies of $\F_2$ and
    numbers inside the circle indicate the $s$-degree, the Cartan--Eilenberg Ext-degree.
    The red dot indicates a class that is possibly zero.
    Middle: 
    Each solid square $\blacksquare$ denotes a copy of $\F_2[\tau]$. 
    Hollowed squares $\square$ denotes possible copies of $\F_2[\tau]$.}
  \label{fig:3charts}    
\end{figure}

% \begin{itemize}
%     \item
%       The horizontal degree is the topological stem $t-a$, shifted by $3\cdot 2^j$. 
%     \item
%       The vertical degree is the Adams filtration $a$. 
%     \item Left and Right: Each dot $\bullet$ denotes a copy of $\F_2$. 
%     \item Middle: 
%       \begin{itemize}
%       \item Each solid square $\blacksquare$ denotes a copy of $\F_2[\tau]$. 
%       \item Hollowed square $\square$ denotes possible copies of $\F_2[\tau]$. 
%       \item Circle $\circ$ and double circle $\circledcirc$ denote possible copies of $\F_2$ and $\F_2[\tau]/\tau^2$,\\ with numbers inside indicating their $s$-degrees.
%       \end{itemize}
%     \end{itemize}

% \begin{rmk}
%   In fact, using arguments that are similar to the proof of Theorem~\ref{E2 descriptions}~$(2)$~$(c)$, we can further rule out a few more possible groups in the statement of Theorem~\ref{E2 descriptions}~$(2)$, such as the possible copies of $\F_2$ in degree $(a, \ t-a, \ s, \ w) = (4, \ 3\cdot 2^j - 3, \ 3, \ 3\cdot 2^{j-1})$ of Theorem~\ref{E2 descriptions}~$(2)$~$(a)$. But for the purpose of proving our main Theorem~\ref{thm: d4 hj3}, we do not need such arguments. 
% \end{rmk}

As an amplification of the first four lines of the table in \Cref{E2 descriptions}(2) we have the following corollary.

\begin{cor} \label{cor:E2 inj}
  The Betti realization map from 
  the $E_2$-page of the motivic Adams sseq for $S^{0,0}$ to
  the $E_2$-page of the Adams sseq for $S^{0}$ is
  injective in the tridegrees of $h_0g_{j-2}$, $h_0^2 g_{j-2}$, $h_0^3 g_{j-2}$ and their $\tau$-multiples.
\end{cor}


%% In order to proceed we need two more inputs from the later sections of this paper.
%% In \Cref{sec:algN} we will prove \Cref{diff motivic ctau} which computes 
%% prove Theorem~\ref{thm: d4 hj3}, we also need the following theorems.

% section 7 input
\begin{thm} \label{diff motivic ctau}
  Let $j \geq 6$.
  In the motivic Adams sseq for $S^{0,0}/\tau$, 
  \begin{enumerate}
  \item $d_2(h_j^3) = 0$,
  \item $d_3(h_j^3) = 0$ and 
  \item $d_4(h_j^3) = h_0^3 g_{j-2}$.
  \end{enumerate}
\end{thm}

% section 8 input
\begin{thm} \label{prop:diff classical}
  Let $j \geq 6$. In the classical Adams sseq for $S^{0}$, 
  \begin{enumerate}
  \item $d_2(h_j^3) = 0$,
  \item $d_3(h_j^3)$ is either $0$ or $h_0^2g_{j-2}$ and
  \item $d_4(h_j^3)$ is either $0$ or $h_0^3g_{j-2}$ (if defined).
  \end{enumerate}  
\end{thm}

With all our inputs ready we now begin proving our main theorem.

% alg N entering diffs
\begin{lem} \label{lem:ctau-entering-diffls}
  Let $j \geq 6$.
  In the motivic Adams sseq for $S^{0,0}/\tau$, 
  there are no non-zero $d_2$ or $d_3$-differentials entering the tridegrees of
  $h_0g_{j-2}$, $h_0^2g_{j-2}$ or $h_0^3 g_{j-2}$.
\end{lem}

\begin{proof}
  Using \Cref{E2 descriptions}(3) we can read off that the only potential sources for a $d_2$ or $d_3$-differential entering the tridegree of one of $h_0g_{j-2}$, $h_0^2g_{j-2}$ or $h_0^3 g_{j-2}$ are
  $h_j^3$, $h_0h_j^3$ and $h_0^2h_j^3$.
  The lemma now follows as a corollary of \Cref{diff motivic ctau}(1,2).
\end{proof}



% d2 differentials
\begin{lem} \label{lem:d2 diffls}
  Let $j \geq 6$.
  \begin{enumerate}
  \item In the motivic Adams sseq for $S^{0,0}$, $d_2(h_j^3) = 0$.
  \item For $j \neq 7$, in the Adams sseq for $S^0$, there are no non-zero $d_2$-differentials entering the bidegrees of $h_0g_{j-2}$, $h_0^2 g_{j-2}$ and $h_0^3 g_{j-2}$.
  \item For $j \neq 7$, in the motivic Adams sseq for $S^{0,0}$, there are no non-zero $d_2$-differentials entering the tridegrees of $h_0g_{j-2}$, $h_0^2 g_{j-2}$, $h_0^3 g_{j-2}$ or their $\tau$-multiples.
  \item The Betti realization map from
    the $E_3$-page of the motivic Adams sseq for $S^{0,0}$ to
    the $E_3$-page of the Adams sseq for $S^{0}$ is
    injective in the tridegrees of $h_0g_{j-2}$, $h_0^2 g_{j-2}$, $h_0^3 g_{j-2}$ and their $\tau$-multiples.
  \end{enumerate}  
\end{lem}

\begin{proof}
  Using \Cref{cor:E2 inj} we can deduce (1) from \Cref{prop:diff classical}(1).  
  For $j \neq 7$  we can read off from \Cref{E2 descriptions}(1) that the only potential sources for Adams differential entering the bidegrees of $h_0g_{j-2}$, $h_0^2 g_{j-2}$ and $h_0^3 g_{j-2}$ are
  the classes $h_j^3$, $h_0h_j^3$, $h_0^2h_j^3$ and $h_1g_{j-2}$.
  Thus, (2) follows from \Cref{prop:diff classical}(1) and \Cref{exm:Bruner formula} (which tells us that $d_2(g_{j-2}) = 0$).
  The injectivity from \Cref{cor:E2 inj} implies that any entering motivic $d_2$-differential would induce an entering classical $d_2$-differential, therefore (2) implies (3).  
  For $j \neq 7$, (4) is obtained by combining (2), (3) and \Cref{cor:E2 inj}.

  For the $j=7$ case of (4) we observe that the additional generator $h_8D_3(1)$ has degree
  $(a,t,w) = (5, 3 \cdot 2^7 + 2, 3 \cdot 2^6 + 1)$ and that
  in order for this differential to create $\tau$-torsion on the motivic Adams $E_3$-page its target must be $\tau$-divisible.
  This would force us to have a non-trivial class in degree
  $(a,t,w) = (7, 3 \cdot 2^7 + 3, 3 \cdot 2^6 + 2)$.
  On the other hand, such a class would contradict the fact that
  the motivic Adams $E_2$-page vanishes when $t < 2w$.
\end{proof}

% d3 differentials
\begin{lem} \label{lem:d3 diffls}
  Let $j \geq 6$.
  \begin{enumerate}
  \item In the classical Adams sseq for $S^{0}$, $d_3(h_j^3) = 0$.
  \item In the motivic Adams sseq for $S^{0,0}$, $d_3(h_j^3) = 0$.
  \item In the classical Adams sseq for $S^0$, there are no non-zero $d_3$-differentials entering the bidegrees of $h_0^2 g_{j-2}$ and $h_0^3 g_{j-2}$.
  \item In the motivic Adams sseq for $S^{0,0}$, there are no non-zero $d_3$-differentials entering the tridegrees of $h_0^2 g_{j-2}$, $h_0^3 g_{j-2}$ or their $\tau$-multiples.
  \item The Betti realization map from
    the $E_4$-page of the motivic Adams sseq for $S^{0,0}$ to
    the $E_4$-page of the Adams sseq for $S^{0}$ is
    injective in the tridegrees of $h_0g_{j-2}$, $h_0^2 g_{j-2}$, $h_0^3 g_{j-2}$ and their $\tau$-multiples.
  \end{enumerate}  
\end{lem}

\begin{proof}
  We begin with (2).
  The injectivity statement from \Cref{lem:d2 diffls}(4) allows us to upgrade
  \Cref{prop:diff classical} to the claim that
  in the motivic Adams sseq for $S^{0,0}$, $d_3(h_j^3)$ is either $0$ or $h_0^2g_{j-2}$.
  Next we examine the reduction mod $\tau$ map to the motivic Adams sseq for $S^{0,0}/\tau$.  
  From \Cref{E2 descriptions}(3) and \Cref{lem:ctau-entering-diffls}
  we know that $h_0^2g_{j-2}$ is non-zero
  on the $E_3$-page of the motivic Adams sseq for $S^{0,0}/\tau$.
  Thus, (2) follows from \Cref{diff motivic ctau}(2) which tells us that
  $d_3(h_j^3) = 0$ in the motivic Adams sseq for $S^{0,0}/\tau$.
  (1) follows from (2) by Betti realization.
        
  From \Cref{E2 descriptions}(1) we can read off that the only potential sources for Adams differential entering the bidegrees of $h_0g_{j-2}$, $h_0^2 g_{j-2}$ and $h_0^3 g_{j-2}$ are
  the classes $h_j^3$ and $h_0h_j^3$.
  Thus, (3) follows from (1).
  The injectivity from \Cref{lem:d2 diffls}(4) implies that any entering motivic $d_3$-differential would induce an entering classical $d_3$-differential, therefore (3) implies (4).  
  (5) is obtained by combining (3), (4) and \Cref{lem:d2 diffls}(4).
\end{proof}

\begin{thm} \label{thm:motivic d4}
  Let $j \geq 6$.
  In the motivic Adams sseq for $S^{0,0}$ we have
  \[ d_4(h_j^3) = h_0^3g_{j-2} \neq 0 \]
  which Betti realizes to the non-trivial Adams differential of \Cref{thm: d4 hj3}
\end{thm}

\begin{proof}
  From Lemmas \ref{lem:d2 diffls}(1) and \ref{lem:d3 diffls}(2) we know that the class $h_j^3$ survives to the $E_4$-page of the motivic Adams sseq.
  Under the reduction mod $\tau$ map to $E_4$-page of the motivic Adams sseq for $S^{0,0}/\tau$
  the class $d_4(h_j^3)$ is sent to a non-zero class by \Cref{diff motivic ctau}(3) and \Cref{lem:ctau-entering-diffls}, therefore $d_4(h_j^3) \neq 0$.
  On the other hand, the injectivity statement from \Cref{lem:d3 diffls}(5) allows us to upgrade
  \Cref{prop:diff classical}(3) to the claim that
  in the motivic Adams sseq for the sphere $d_4(h_j^3)$ is either $0$ or $h_0^3g_{j-2}$.
  It follows that
  \[ d_4(h_j^3) = h_0^3g_{j-2} \neq 0. \]
  
  Using the injectivity from \Cref{lem:d3 diffls}(5) again we obtain the desired non-trivial classical Adams $d_4$-differential from the motivic one.
  \qedhere
  
  % Using Theorem~\ref{diff motivic sphere}, the element $h_j^3$ survives to the $E_4$-page of the motivic Adams spectral sequence, in particular we obtain a class $d_4(h_j^3)$ on the $E_4$-page of the motivic Adams spectral sequence. Now consider the map between motivic Adams spectral sequences induced by the quotient map $S^{0,0} \rightarrow S^{0,0}/\tau$, using Theorem~\ref{diff motivic ctau} we can read off that
  % $d_4(h_j^3)$ maps to $h_0^3g_{j-2}$ under this reduction map.
  % Now we apply \Cref{lem:zig-zag-d4} in order to conclude that the Betti realization of $d_4(h_j^3)$ is non-zero. As a consequence we learn that in the classical Adams spectral sequence we have $d_4(h_j^3) \neq 0$.

  % To complete the proof we now need to identify the target of this classical Adams $d_4$-differential.
  % By Theorem~\ref{diff classical}~$(2)$, we have $h_0^2$ divides $d_4(h_j^3)$. From the description of the classical $\Ext_{\A}$ in Theorem~\ref{E2 descriptions}~$(1)$, the only non-zero class that is divisible by $h_0^2$ is $h_0^3g_{j-2}$. In other words, $d_4(h_j^3) = h_0^3g_{j-2}$.
\end{proof}

\begin{rmk}
One may prove directly that the element $h_0^3g_{j-2}$ survives to the classical Adams $E_4$-page, but it is not logically necessary for the proof of \Cref{thm:motivic d4}---\Cref{thm:motivic d4} in fact implies that $h_0^3g_{j-2}$ does not support a nonzero $d_2$ or $d_3$-differential. %If , then it would contradict to conclusion of Theorem~\ref{thm: d4 hj3} that $h_0^3g_{j-2}$ is a permanent cycle and being nonzero on the $E_4$-page.
\end{rmk}



%% In the rest of this paper, we prove Theorems~\ref{E2 descriptions}, \ref{diff classical}, \ref{diff motivic sphere}, \ref{diff motivic ctau}. In Section~\ref{sec:cess}, we set up the Cartan-Eilenberg spectral sequence and prove part $(1)$ of Theorem~\ref{E2 descriptions}. We also prove Theorem~\ref{diff motivic ctau}~$(3)$ at the end of Section~\ref{sec:cess}. In Section~\ref{sec:mcess}, we set up the motivic analog of the Cartan-Eilenberg spectral sequence and prove parts $(2)(3)$ of Theorem~\ref{E2 descriptions}. In Section~\ref{sec:e4}, we prove Theorem~\ref{diff classical}~$(3)$ and we use Theorem~\ref{diff motivic ctau}~$(1)$ and the motivic zigzag strategy to prove Theorem~\ref{diff classical}~$(1)$ and Theorem~\ref{diff motivic sphere}. These proofs are similar to but easier than the proof of our main Theorem~\ref{thm: d4 hj3}. What remains to be proved are Theorem~\ref{diff motivic ctau}~$(1)(2)$ and Theorem~\ref{diff classical}~$(2)$. In Section~\ref{sec:alg-nov-diff}, we study the key algebraic Novikov differential through the cobar complex, and prove Theorem~\ref{diff motivic ctau}~$(1)(2)$. In Section~\ref{sec:kervaire}, we revisit Barratt--Jones--Mahowald's inductive approach on the Kervaire invariant classes and prove Theorem~\ref{diff classical}~$(2)$.




%% --------------end re-arranged part---------------------------------------


% \begin{thm} \label{diff classical}
% Suppose that $j \geq 6$. In the classical Adams spectral sequence for $S^{0}$, 
% \begin{enumerate}
% \item $d_2(h_j^3) = 0$ and $d_3(h_j^3) = 0$,
% \item $d_4(h_j^3)$ is divisible by $h_0^2$.
% \item There are no non-zero $d_2$ or $d_3$-differentials entering the bidegree of $h_0^3 g_{j-2}$.\todo{This needs to get proved.}
% \end{enumerate}
% \end{thm}

% \begin{thm} \label{diff motivic sphere}
% Suppose that $j \geq 6$. In the motivic Adams spectral sequence for $S^{0,0}$, 
% %\begin{enumerate}
% %\item 
% $$d_2(h_j^3) = 0 \ \text{and} \ d_3(h_j^3) = 0.$$
% %\item $d_4(h_j^3)$ is not the target of any $d_2$ or $d_3$-differential.
% %\end{enumerate}
% \end{thm}


% Now we use Theorems~\ref{E2 descriptions}, \ref{diff classical}, \ref{diff motivic sphere}, \ref{diff motivic ctau} to prove Theorem~\ref{thm: d4 hj3}.



% \begin{lem} \label{lem:zig-zag-d4} 
%   Suppose $y$ is a class on the $E_4$-page of the motivic Adams spectral sequence for $S^{0,0}$ in degree $(a,t-a, w) = (7, 3 \cdot 2^j -4, 3 \cdot 2^{j-1})$
%   that maps to $h_0^3g_{j-2}$ under the reduction mod $\tau$ map to the $E_4$-page of the motivic Adams spectral sequence for $S^{0,0}/\tau$.
%   Then, the Betti realization of $y$ is a non-zero class on the $E_4$-page of the classical Adams spectral sequence.
% \end{lem}

% \begin{proof}
%   Let ${}^{\mathrm{mA}}Z_4$ (resp. ${}^{\mathrm{A}}Z_4$ and ${}^{\mathrm{aN}}Z_4$) denote the subgroup of classes on the $E_2$-page of the motivic Adams spectral sequence for $S^{0,0}$ (resp. classical Adams and motivic Adams for $S^{0,0}/\tau$) in degree $(a,t-a, w) = (7, 3 \cdot 2^j -4, 3 \cdot 2^{j-1})$ on which the $d_2$ and $d_3$ differential vanish. Then we may consider the following diagram
%   \[ \begin{tikzcd}
%       {}^{\mathrm{aN}}E_2 & {}^{\mathrm{mA}}E_2 \ar[l] \ar[r] & {}^{\mathrm{A}}E_2 \\
%       {}^{\mathrm{aN}}Z_4 \ar[u, hook] \ar[d, two heads, "\cong"] & {}^{\mathrm{mA}}Z_4 \ar[l] \ar[r] \ar[u, hook] \ar[d, two heads] & {}^{\mathrm{A}}Z_4 \ar[u, hook] \ar[d, two heads, "\cong"] \\
%       {}^{\mathrm{aN}}E_4 & {}^{\mathrm{mA}}E_4 \ar[l] \ar[r] & {}^{\mathrm{A}}E_4
%     \end{tikzcd} \]
%   where the bottom left vertical map is an equivalence by \Cref{diff motivic ctau}(3)
%   and the bottom right vertical map is an eqvuialence by \Cref{diff classical}(3).

%   Choose a class $\widetilde{y}$ in ${}^{\mathrm{mA}}Z_4$ which maps to $y$ in ${}^{\mathrm{mA}}E_4$.
%   It will now suffice for us to argue that the image of $\widetilde{y}$ in ${}^{\mathrm{aN}}Z_4$ is non-zero. In fact, using the vertical inclusions it suffices work at the level of the $E_2$-page. 
%   Theorem~\ref{E2 descriptions}(2d) lets us decompose ${}^{\mathrm{mA}}E_2$ as an $\F_2[\tau]$-module as follows
%   \[ \Ext^{7,\,3\cdot 2^j - 4,\,*}_{A^\textup{mot}}(\mathbb{F}_2[\tau], \mathbb{F}_2[\tau]) \cong  \F_2[\tau]\{h_0^3g_{j-2}\} \oplus M. \]
%   On $E_2$-pages the kernel of the reduction mod $\tau$ map
%   \[ E_2^{a,t,w}(S^{0,0}) \to E_2^{a,t,w}(S^{0,0}/\tau) \]
%   is exactly the classes which are divisible by $\tau$, therefore we have $\widetilde{y} = h_0^3g_{j-2} + m$ for some $m \in M$.
%   Finally, the splitting of $\F_2[\tau]$-modules above lets us conclude that the $\widetilde{y}$ remains non-zero upon inverting $\tau$ since $h_0^3g_{j-2}$ remains non-zero upon inverting $\tau$.
%   \qedhere
  
%   % on the $E_2$-page of the motivic Adams spectral sequence for $S^{0,0}$ such that $d_2(\widetilde{y})=0$, $d_3(\widetilde{y})=0$ and $\widetilde{y}$ maps to $y$ on the $E_4$-page.
     
%   % Choose a class $\widetilde{y}$ on the $E_2$-page of the motivic Adams spectral sequence for $S^{0,0}$ such that $d_2(\widetilde{y})=0$, $d_3(\widetilde{y})=0$ and $\widetilde{y}$ maps to $y$ on the $E_4$-page.
%   % As the image of 

%   % Since $d_4(h_j^3)$ is in $(a, \ t-a)$-degree $(7, \ 3 \cdot 2^j -4)$,
%   % we apply Theorem~\ref{E2 descriptions}~(2)~(d).
%   % As a module over $\F_2[\tau]$, we may write 
%   % $$\Ext^{7,\ 3\cdot 2^j - 4, \ *}_{A^\textup{mot}}(\mathbb{F}_2[\tau], \mathbb{F}_2[\tau]) = F \oplus T.$$
%   % Here $F$ is free over $\F_2[\tau]$, generated by $h_0^3 g_{j-2}$.
% \end{proof}


% \begin{proof}[Proof of Theorem~\ref{thm: d4 hj3}] 
% First consider the motivic Adams spectral sequence for $S^{0,0}/\tau$. By Theorem~\ref{diff motivic ctau}$(1)$, $h_j^3$ survives to the $E_4$-page. By Theorem~\ref{diff motivic ctau}~$(2)$ and $(3)$, $h_0^3 g_{j-2}$ is the target of a differential, and therefore is a permanent cycle. So in particular, it is a $d_2$ and $d_3$-cycle. Moreover, $h_0^3 g_{j-2}$ is not the target of any $d_2$ or $d_3$-differential. Therefore, $h_0^3 g_{j-2}$ survives to the $E_4$-page. This tells us
% $$d_4(h_j^3) = h_0^3 g_{j-2} \neq 0,$$
% in the motivic Adams spectral sequence for $S^{0,0}/\tau$.

% Next consider the motivic Adams spectral sequence for $S^{0,0}$. By Theorem~\ref{diff motivic sphere}, the element $h_j^3$ survives to the $E_4$-page. Now consider the map between motivic Adams spectral sequences induced by the quotient map $S^{0,0} \rightarrow S^{0,0}/\tau$, we have that 
% $$d_4(h_j^3) \neq 0,$$
% in the motivic Adams spectral sequence for $S^{0,0}$.

% Since $d_4(h_j^3)$ is in $(a, \ t-a)$-degree $(7, \ 3 \cdot 2^j -4)$, we apply Theorem~\ref{E2 descriptions}~(2)~(d). As a module over $\F_2[\tau]$, we may write 
% $$\Ext^{7,\ 3\cdot 2^j - 4, \ *}_{A^\textup{mot}}(\mathbb{F}_2[\tau], \mathbb{F}_2[\tau]) = F \oplus T.$$
% Here $F$ is free over $\F_2[\tau]$, generated by $h_0^3 g_{j-2}$, and elements $f_{i,l}, \ i = 0,\ 2, \ 4, \ 6, \ l \in C_i$ for some index sets $C_i$, and $T$ is $\tau$-torsion over $\F_2[\tau]$, generated by elements $a_{i,l}, \ i = 4,\ 6,\  l \in A_i$ and $b_{6, l}, \ l \in B_6$ for some index sets $A_4, \ A_6, \ B_6$, where $\tau \cdot b_{6,l} \neq 0$ (So $a_{i,l}$ generates $\F_2$ and $b_{6,l}$ generates $\F_2[\tau]/\tau^2$). The $(s,\ w)$-degrees of $f_{i,l}, \ a_{i,l}$ and $b_{i,l}$ are $(i,\ 3\cdot 2^{j-1}-2+\frac{i}{2})$. It is clear that if we further concentrate on weight $3\cdot 2^{j-1}$, the elements 
% $$h_0^3 g_{j-2}, \ f_{4,l}, \ l \in C_4, \ \tau \cdot f_{6,l}, \ l \in C_6, \ a_{4,l}, \ l \in A_4, \ \tau\cdot b_{6,l}, \ l \in B_6$$ 
% form a basis for $\Ext^{7,\ 3\cdot 2^j - 4, \ 3\cdot 2^{j-1}}_{A^\textup{mot}}(\mathbb{F}_2[\tau], \mathbb{F}_2[\tau])$ over $\F_2$.

% Since motivic Adams differentials preserve weight, we know $d_4(h_j^3)$ has the same weight as $h_j^3$, i.e., $3\cdot 2^{j-1}$. Therefore, we have $d_4(h_j^3)$ as a linear combination of of the basis above.
% $$d_4(h_j^3)= \alpha \cdot h_0^3 g_{j-2} + f_4 + \tau \cdot f_6 + a_4 + \tau \cdot b_6,$$
% where $\alpha \in \{0,\ 1\}$ to be determined, $f_4, \ f_6,\ a_4,\ b_6$ are linear combinations of elements of $f_{4,l}, \ l \in C_4, \ f_{6,l}, \ l \in C_6, \ a_{4,l}, \ l \in A_4, \ b_{6,l}, \ l \in B_6$ respectively. Note that since the expression of $d_4(h_j^3)$ is a nonzero element in the $E_4$-page, it does not support nonzero $d_2$ or $d_3$-differentials, and it is hit by nonzero $d_2$ and $d_3$-differentials.

% Map the expression for $d_4(h_j^3)$ from the $E_2$-page of the motivic Adams spectral sequence of $S^{0,0}$ to that of $S^{0,0}/\tau$. The elements $h_0^3 g_{j-2}, \ f_4, a_4$ map nontrivially to their corresponding elements and $\tau \cdot f_6,\ \tau \cdot b_6$ map to zero, since they are $\tau$-multiples. Because $d_4(h_j^3) = h_0^3 g_{j-2}$ in the motivic Adams spectral sequence of $S^{0,0}/\tau$, we must have $\alpha = 1$ and 
% $$d_4(h_j^3)= h_0^3 g_{j-2} + \tau \cdot f_6 + \tau \cdot b_6,$$
% in the motivic Adams spectral sequence of $S^{0,0}$.

% At last, we consider the classical Adams spectral sequence for $S^{0}$. By Theorem~\ref{diff classical}~$(1)$, the element $h_j^3$ survives to the $E_4$-page. Consider the map of Adams spectral sequences from the motivic one for $S^{0,0}$ to the classical one for $S^{0}$ through the Betti realization. We have that 
% $$d_4(h_j^3) = h_0^3g_{j-2} + f_6 \neq 0.$$
% By Theorem~\ref{diff classical}~$(2)$, we have $h_0^2$ divides $d_4(h_j^3)$. From the description of the classical $\Ext$ in Theorem~\ref{E2 descriptions}~$(1)$, the only possibility for $d_4(h_j^3)$ being nonzero and divisible by $h_0^2$ is that $d_4(h_j^3) = h_0^3g_{j-2}$. In other words, $f_6 =0$. This completes the proof. 
% \end{proof}





